
\chapter{Related Work}
\label{ch:related}

\instructions{Page budget for Related Work: 10 - 15 pages
%
\begin{itemize}
    \item Discuss the relevant literature related with your problem, in a coherent way to demonstrate that you have given due consideration to know: your problem and the possible sub-problems of interest in this work, the previous approaches, and their benefits and missing points or drawbacks.
    \item It is not about concatenating summaries of papers. Rather, it is about to discuss the problem, its details and challenges (difficulties and/or opportunities) you focus on, all this by relying on previous work, comparing what some did, what some else improved or added, what is missing, and so on.
    \item Moreover, mind to keep the focus in the areas you care during your discussion (since you control the story), and to remark what you consider useful, interesting, or challenging.
    \item In some cases it could be necessary and/or natural to have, first, an introduction (which could discuss works on more general aspects of the problem), and then split the rest of the chapter in sections, each discussing a particular aspect of the problem.
    \item If you define a new task, or create or improve a method, or develop a test collection, or perform a missing major evaluation, explain why it is interesting, why is different or helpful versus previous works.
    \item Mind that the reader may have never heard about these things. You need to discuss them in such detail that it is possible to follow later parts of the thesis without having to consult external resources.
    \item Always use natbib!
    \item Use \texttt{\textbackslash\.citet\{\}} for textual citation. For example, ``\citet{Balog:2018:Book} proposed...''
    \item Use \texttt{\textbackslash\.citep\{\}} for parenthetical citation. For example, ``In \citep{Zhang:2020:KDD} the idea of ..''
    \item Here is an example of a PhD thesis:  \citet{Maxwell:2019:PhDThesis} 
    \item Here is how a Journal article would look in the Bibliography: \citet{Sanderson:2010:FnTIR}
    \item Never write out Smith et al., there is a \texttt{\textbackslash\.citeauthor{}\{\}} command for that (but most likely what you're looking for is actually \texttt{\textbackslash\.citet\{\}}.
    \item Check the corresponding Bib entry, for citing online documentation~\citep{Rasa:2022:doc}; if you just need to include an URL of a website/code/dataset, use a footnote instead.\footnote{\url{https://rasa.com}}
\end{itemize}
}

The rapid growth of social media platforms and online content-sharing systems has significantly increased the spread of misinformation. False or misleading information can now propagate globally within minutes, often reaching large audiences before verification occurs. This problem has become particularly critical in domains such as public health, politics, crisis reporting, and social events, where misinformation can influence public opinion, create panic, or distort factual understanding.
A major challenge in modern misinformation is that it is increasingly multimodal. Many misleading claims are accompanied by images, videos, or screenshots that provide apparent visual credibility to false narratives. In numerous cases, authentic images are reused in misleading contexts, attached to fabricated stories, or combined with inaccurate textual claims. As a result, verifying such claims requires not only textual reasoning but also understanding the relationship between visual and textual information.
Traditional manual fact-checking performed by professional organizations is effective but difficult to scale due to the enormous volume of online content. This limitation has motivated extensive research into automated fact-checking systems. Early automated approaches primarily focused on textual verification, where the goal was to retrieve supporting evidence and determine whether a claim was true or false. However, text-only systems often struggle when misinformation relies heavily on visual manipulation, image miscontextualization, or cross-modal inconsistencies.
Recent advances in large language models, retrieval systems, and vision-language models have enabled the development of more sophisticated multimodal fact-checking pipelines. In particular, retrieval-based approaches have become increasingly important because the quality of retrieved evidence directly influences downstream verification performance. Modern systems therefore combine semantic retrieval, lexical retrieval, and multimodal retrieval techniques to improve evidence coverage and robustness.
Despite these advances, several challenges remain unresolved. Existing systems often:
rely heavily on textual retrieval while underutilizing visual evidence,
struggle to identify reused or miscontextualized images,
retrieve redundant or noisy evidence,
and fail to effectively combine semantic, lexical, and visual retrieval signals.
Furthermore, many image-text retrieval models are optimized for short captions rather than the longer contextual evidence commonly encountered in fact-checking datasets.
This thesis addresses these challenges by proposing a hybrid multimodal retrieval framework for automated fact-checking. The proposed approach integrates dense retrieval using BGE, sparse retrieval using SPLADE, and image-text retrieval using LongCLIP to retrieve diverse and relevant evidence from both textual and visual sources. The system further combines retrieval outputs using Reciprocal Rank Fusion (RRF), cross-encoder reranking, and Maximum Marginal Relevance (MMR) to improve retrieval quality and evidence diversity.
The purpose of this chapter is to review the existing literature relevant to automated fact-checking, multimodal misinformation analysis, retrieval-augmented systems, and multimodal retrieval techniques. The discussion focuses on the strengths and limitations of previous approaches, the challenges associated with multimodal fact verification, and the research gaps that motivate the proposed system.

\section{Traditional Fact-Checking}

Traditional fact-checking is the process of manually verifying the accuracy of a claim.Before the emergence of automated fact-checking approaches, verification of information was primarily performed manually by journalists, researchers, and professional fact-checking organizations. Traditional fact-checking systems rely heavily on human expertise to investigate claims, analyze evidence, consult reliable sources, and publish verification reports. These systems play an important role in combating misinformation by providing high-quality, carefully reviewed assessments of public claims and viral content.

Manual fact-checking, while serving as a cornerstone of journalistic accountability, is increasingly hindered by structural and temporal constraints that limit its effectiveness in the modern digital landscape. The following analysis outlines the primary limitations identified in the sources.

\begin{itemize}
    \item \textbf{Temporal Latency and the "Speed Gap"}: The most significant limitation of the manual fact-checking process is its inherent lack of speed. Fact-checking is described as an intellectually demanding and laborious task that requires extensive research and advanced writing . A typical check requires approximately one full day to research and write, with more complex claims demanding two or more days . This creates a critical "speed gap" between the moment a false statement is made and the publication of a correction . During this interval, misinformation can circulate unchecked, allowing false narratives to take root before they can be officially refuted .
    
    \item \textbf{Constraints on Scalability and Volume}: Manual fact-checking is not scalable enough to match the volume of misinformation. Human journalists are unable to monitor and verify the sheer quantity of claims produced during live events, such as debates or constant social media updates, in real-time . Because the process is so time-intensive, a single news organization can only produce a limited number of checks per day, which is insufficient for the high velocity and scale of misinformation today .
    
    \item \textbf{Lack of Structural Data and Centralization}: The efficacy of manual fact-checking is further diminished by fragmented dissemination. Most traditional fact-checkers utilize older, blog-style platforms that publish work as unstructured plain text . This format lacks the "structured journalism" metadata—such as specific tags for the speaker, statement, and evidence—necessary for automated tools to easily digest and reuse the information .
Furthermore, there is no single central repository for these verified facts . Because checks are scattered across various individual news archives, they are difficult for the public to find during live events and lead to an inefficient duplication of effort, where different organizations may spend hours independently verifying the same claim because there is no shared database .
    
    \item \textbf{Structural Constraints and Limited Reach}: From a structural perspective, traditional fact-checking lacks the direct line of communication to users that social media platforms enjoy. Because they are centralized and often detached from the networks where misinformation spreads, these organizations have less reach and no consistent, data-driven measure of their actual impact on public discourse. Unlike decentralized models that can leverage network data to ensure a diversity of perspectives, traditional models rely on a limited number of employees, which restricts their ability to scale and maintain perceived ideological neutrality ~\cite{QuestToAutomateFactChecking2015Oct}.
    \item \textbf{Ideological Capture and Perceptions of Bias}: Traditional fact-checking organizations face intense scrutiny regarding perceived and actual political bias. This demographic reality leads to "ideological capture," where the lack of diverse viewpoints within a small, centralized staff—the "bottleneck" of analysis—contributes to a public perception that fact-checkers are partisan actors rather than neutral arbiters. Consequently, these institutions lack effective internal mechanisms to combat such perceptions, which ultimately reduces their ability to combat misinformation effectively ~\cite{FactChecking2025Mar}.
\end{itemize}

To overcome these limitations, there has been a growing interest and a rapid development in automated fact-checking systems that can operate at scale and in real-time. These systems leverage advances in natural language processing, machine learning, and information retrieval to assist human fact-checkers or to perform verification tasks autonomously. 

\section{Automated Fact-Checking}

Automated fact-checking is the pursuit of a computer-based platform capable of detecting factual claims as they appear in real-time and providing an immediate accuracy rating without human intervention. Often referred to in the research community as the "Holy Grail" of computational journalism, a fully automated system is characterized by being instant, accurate, and accountable—meaning it self-documents its data sources and analysis for transparency ~\cite{QuestToAutomateFactChecking2015Oct}.

Automated fact checking is a challenging task that requires the system to perform several complex subtasks, including claim detection, evidence retrieval, and claim verification. The goal is to determine the veracity of a claim by retrieving relevant evidence from a large corpus of documents and analyzing the relationship between the claim and the evidence.

The "ever-increasing amounts of textual information" available today, combined with the ease of sharing it online, has created a high demand for automated verification. 

Vlachos and Riedel (2014) constructed a dataset for claim verification consisting of 106 claims, selecting data from fact-checking websites such as PolitiFact, taking advantage of the labelled claims available there. However, in order to de-velop claim verification components we typically require the justification for each verdict, includ-ing the sources used. While this information is usually available in justifications provided by the journalists, they are not in a machine-readable form ~\cite{Vlachos2014Jun}.

Wang (2017) extended this approach by including all 12.8K claims available by Politifact via its API, however the justification and the evidence contained in it was ignored in the experiments as it was not machine-readable. Instead, the claims were classified considering only the text and the metadata related to the person making the claim. While this rendered the task amenable to current NLP/ML methods, it does not allow for verification against any sources and no evidence needs to be returned to justify the verdicts ~\cite{Wang2017Jul}.

The Fake News challenge (Pomerleau and Rao, 2017) modelled verification as stance classification: given a claim and an article, predict whether the article supports, refutes, observes (neutrally
states the claim) or is irrelevant to the claim. It consists of 50K labelled claim-article pairs, combining 300 claims with 2,582 articles. The claims and the articles were curated and labeled by journalists and the dataset was first proposed
by Ferreira and Vlachos (2016), who only classified the claim w.r.t. the article headline instead of the whole article and the systems were provided with the sources to verify against, instead of having to retrieve them ~\cite{PomerleauDelip2017Jun}.

~\citet{Thorne2018Jun} introduced a new style of fact-checking characterized by its scale and retrieval-based complexity FEVER: Fact Extraction and VERification. It consists of 185,445 claims verified against Wikipedia ~\cite{Thorne2018Jun}.

Key components of the FEVER-style approach include:

\begin{itemize}
    \item \textbf{Three-Way Classification}: Claims are labeled as SUPPORTED, REFUTED, or NOTENOUGHINFO.
    
    \item \textbf{Evidence Retrieval}: systems are not given the evidence; they must retrieve it from a massive document collection.
    \item \textbf{Sentence-Level Justification}: For a verdict to be considered correct, the system must return the specific sentences that form the necessary evidence.
    \item \textbf{Complex Evidence Composition}: FEVER-style claims often require multi-hop inference, where evidence must be composed from multiple sentences (16.82\% of cases) or even multiple pages (12.15\% of cases).

\end{itemize}

By requiring both information retrieval and textual inference, FEVER serves as a challenging testbed for developing new methods in reading comprehension and claim extraction.

However this dataset contain purpose-made claims derived from Wikipedia, and thus are not representative of the real-world misinformation that circulates online. The claims in FEVER are artificially generated by altering sentences from Wikipedia, which may not capture the complexity, nuance, and multimodal nature of real-world misinformation.

~\citet{Schlichtkrull2023May} introduced a new dataset for fact checking, AVERITEC, which combines real-world claims with realistic evidence retrieved from the web, as well as justifications for veracity labels. We formulate retrieval as question generation and answering, providing a structured representation of the evidence and reasoning supporting or refuting the claim. The free-text justifications detail how the evidence is used to reach the verdict, including cases of conflicting
evidence, matching best practises for human fact-checkers. This dataset consists of 4,568 real world claims. They have avoided, common pitfalls including  context dependence,
evidence insufficiency, and temporal leakage, which are common in other datasets. The claims in AVERITEC are sourced from real-world fact-checking organizations, making it more representative of the types of misinformation encountered in practice. Additionally, the dataset includes realistic evidence retrieved from the web, which adds to its authenticity and relevance for developing practical fact-checking systems ~\cite{Schlichtkrull2023May}.

These datasets however focus on textual claims and evidence, while in the real world, misinformation often involves multimodal content, such as images, videos, and screenshots. This has led to the emergence of multimodal fact-checking research, which aims to verify claims that involve both textual and visual information.

\section{Multimodal Fact Checking}

Recent studies estimate that approximately 80\% of online claims are multimodal involving both text and media ~\citet{Dufour2024May} and among these media, images are the most prevalent media type. So these days, fact checking cannot be done only using the textual claims and textual evidences, it needs a combined approach involving both image claims, image evidences and image-text alignment.

~\citet{Cao2025May} introduced a new dataset for multimodal fact checking, AVERIMATEC (Automatic Verification of Image-Text Claims) which consists of 1,297 real world image text claims. Each claim
is annotated with question-answer (QA) pairs containing evidence from the web.  In addition, each claim is
annotated with metadata information, a veracity label based on retrieved evidence and a textual
justification, explaining how the verdict is reached.

\section{Retrieval-Augmented Multimodal Fact-Checking Pipelines}

\subsection{Retrieval-Augmented Generation (RAG) for Fact-Checking}

Large Language Models (LLMs) have demonstratred impressive capabilities in natural language understanding and egenration and have been used as fact checking tools.
However according to ~\citet{Proma2025Feb} and ~\citet{Lewis2020May}, Large language models (LLMs) fail in fact-checking primarily because they struggle to ground their responses in real, credible news sources and exhibit a systemic bias in source selection.
When prompted to provide evidence for claims LLMs frequently invent non-existent headlines and articles. This is because LLMs are limited by their training data's "cut-off date."
This means they are often ill-equipped to fact-check recent events that occurred after their training was finalized. LLMs are prone to generating false information with a high degree of confidence.
This becomes dangerous because users tend to place high trust in LLM responses, making them more susceptible to the misinformation the model generates.

This results in the need to external evidences to ground the reasoning of the model, and thus the need for retrieval systems to retrieve relevant evidence from the web or a large corpus of documents. This naturally introduces Retrieval-Augmented Generation (RAG) approaches, where the model retrieves relevant information from an external source and uses it to generate a response. 
Retrieval-Augmented Generation (RAG) is a specialized pipeline designed to enhance the accuracy of large language models (LLMs) by allowing them to access and use external information during the generation process.
One of the major pitfalls of standard LLMs is their tendency to "hallucinate" or invent non-existent sources. RAG helps solve this by integrating external search capabilities, such as web searches, which can provide the model with real, credible evidence to support its claims.
Standard LLMs are limited by "training cut-offs," meaning they cannot accurately fact-check events that occurred after their training data was finalized.
RAG allows the system to retrieve current information from the internet or specific databases, making it more reliable for reporting on recent events.
In summary, RAG acts as a "guardrail" for base systems, moving beyond simple prompting to ensure the model's output is supported by actual external evidence rather than just its internal statistical patterns ~\citep{Lewis2020May}.

In the context of automated fact-checking, retrieval-augmented generation (RAG) has emerged as a powerful approach to enhance the capabilities of large language models (LLMs) by allowing them to access and utilize external information during the generation process. This is particularly important in fact-checking tasks, where the accuracy of the model's output is heavily dependent on the quality and relevance of the retrieved evidence. So RAG
has become a central component in modern fact-checking pipelines, especially in the context of FEVER-style tasks, where the system must retrieve evidence from a large corpus and use it to verify claims.


~\citet{Lewis2020May} used the RAG approach for automated fact checking in the FEVER dataset by ~\citet{Thorne2018Jun}, where RAG was used to retrieve evidences from Wikipedia relating to the claim and then reason over the evidence to classify the claim. For 3-way classification, RAG scores are within 4.3% of
state-of-the-art models, which are complex pipeline systems with domain-specific architectures and
substantial engineering, trained using intermediate retrieval supervision, which RAG does not require. For 2-way classification, we compare against ~\citet{Thorne2018Jun}, who train RoBERTa
to classify the claim as true or false given the gold evidence sentence. RAG achieves an accuracy
within 2.7\% of this model, despite being supplied with only the claim and retrieving its own evidence.
Both the baselines for Fever competition in ~\citet{Schlichtkrull2023May} and ~\citet{Schlichtkrull2024Nov} have used Google Search API to retrieve evidence documents from the internet and have also created retrived evidences from alternative offline knowledge stores, to allow for
inexpensive experimentations with a variety of retrieval strategies, which forms the basis for  transformer models to reason with the evidence to generate a veridct.
Classical TF–IDF/BM25 retrieval limited earlier pipeline QA systems which had a low recall value affecting the retrieval quality. ~\citet{Karpukhin2020Nov} introduced Dense Passage Retrieval (DPR), training dual BERT encoders to embed questions and passages into a shared space. DPR improved top 20 recall by 9–19 pp over BM25 and became the de facto index for RAG models. Modern RAG implementation for evidence retrieval and in fever competitions, vector-based dense retrieval systems were common , along with BM25 ~\citet{Schlichtkrull2024Nov}.


Early work in automated fact-checking primarily focused on textual claims and evidence, but the rise of multimodal misinformation has necessitated the development of systems that can handle both text and images.

Beyond text-only settings, multimodal disinformation research shows that pairing text with images can substantially increase the
perceived credibility and spread of false claims,
motivating fact-checking approaches that incorporate visual evidence  ~\citet{Hameleers2020Feb}.

In multi-modal environments, the intersection of textual claims and visual evidence represents a primary vector for misinformation, where the deception often lies not in the falsehood of a single modality, but in the manipulated contextual connection between them.

\begin{itemize}
    \item \textbf{Image Reuse and Out-of-Context Media}: A prevalent strategy in multi-modal misinformation involves the deployment of out-of-context media, categorized in the literature as a form of fabricated content. In these instances, a news report or social media post utilizes "auxiliary images" that are factually authentic but belong to irrelevant events, effectively repurposing real imagery to support a false current narrative
    
    \item \textbf{Misleading Captions and False Connections}: The category of "False Connection" is identified as one of the most common types of multi-modal misinformation. This occurs when modalities—specifically captions or titles—do not support the accompanying visual content. These misleading captions are frequently designed as clickbait, using provocative or sensationalist headlines to capture user attention despite a lack of contextual consistency with the image. Such discrepancies are often intentionally subtle to exploit the fact that users frequently skim textual content carelessly while being more "favorably attracted" to the visual modality
    
    \item \textbf{Manipulated Narratives and Cross-Modal Discordance}: Sophisticated misinformation often involves manipulated narratives where the individual components (text and image) may be credible in isolation, but their combination creates a deceptive whole. Also manipulated content is created when a valid visual information is edited or altered—such as through Deepfake technology—specifically to support these fabricated narratives. Detecting these narratives requires moving beyond uni-modal analysis to identify cross-modal discordance ~\cite{Abdali2024Nov}.
\end{itemize}

Early work in multimodal fact checking ~\citet{Xu2024Mar} used CLIP, where its visual and text encoders are used to encode the image and the text of the claim to convert to the same feature space and then the attention mechanism and the bilinear pooling are used to fuse the two vectors , and then a classifier is trained on top of this representation to predict the veracity of the claim. However, this approach does not explicitly retrieve evidence from external sources, which limits its ability to provide justifications for its predictions.

similarly the work done by ~\citet{Abdullah2026Jan} introduces an ensemble-based framework that integrates textual and visual data using ViLBERT’s two-stream architecture, incorporates VADER sentiment analysis to detect emotional language, and uses Image–Text Contextual Similarity to identify mismatches between visual and textual elements. These features are processed through the Bi-GRU classifier, Transformer-XL, DistilBERT, and XLNet, combined via a stacked ensemble method with soft voting, culminating in a T5 metaclassifier that predicts the outcome for robustness.
While this approach effectively combines multiple modalities and models, it does not focus on evidence retrieval or the interpretability of its predictions, which are crucial for practical fact-checking applications.

Embedding models like ColPali ~\cite{Faysse2024Jun} encode visual and textual features within multi-vector embeddings
and combine them with the late-interaction mechanism from ColBERT ,allowing for the semantic comparison between texual queries and multimodal image patches and enabling state-of-the-art results on multi-modal retrieval benchmarks.

~\citet{Duwal2025May} propose a graph-based method that evaluates the consistency between the image and the caption by constructing two graph representations: an evidence graph, derived from online textual evidence, and a claim graph, from the claim in the caption. Using graph neural networks (GNNs) to encode and compare these representations, their framework then evaluates the truthfulness of image-caption pairs.

Recent systems combine retrieval with multimodal reasoning and QA generation. in the new AVERIMATEC dataset, ~\citet{Cao2025May} proposed a new approach for multimodal fact checking that explicitly incorporates evidence retrieval. The system consists of three main components: a question generation module that generates questions based on the claim, a question answering module that retrieves evidence from the web to answer the generated questions, and a claim verification module that uses the retrieved evidence to predict the veracity of the claim. This approach allows for more interpretable predictions and provides justifications based on retrieved evidence, addressing some of the limitations of previous methods.

~\citet{Upravitelev2026Mar} have used VLM(Visual Language Model) in their  ADA-AGGR paper to generate a single concise verification question for each claim, and have used an ensemble evidence retrieval comprising of text retreival using BM25 followed by dense retrieval to rerank the top results, a image retrieval for find relevant images based solely on visual data and a textualclaim, the generated question and the claim images are fused into a single query to retriee relevant multimodal evidence from the web.  For Each retrieved image, a VLM is prompted with the image, the original claim and the generated
question to prodeuce a concise answer based strictly on the visual evidence. During this step, the model also performs a binary check; if an image is determined to be irrelevant to the question, it is filtered out to prevent misleading the final verdict.
The system aggregates the top textual and image-based evidence (often limited to 10 items).The VLM analyzes this aggregated evidence to provide a verdict (e.g., Supported, Refuted, or Not Enough Evidence) and a concise justification that cites specific evidence items. To make this process viable for real-world deployment, developers can use aggressive visual downsampling—shrinking evidence images by a factor of up to 8—which can reduce processing time by over 6 times while maintaining comparable verification quality.
However, this approach relies on a single question generation step and may not fully capture the complexity of the evidence or the claim, which can limit its interpretability and robustness.

~\cite{Ullrich2026Mar} follow a similar RAG based approach called AIC CTU, where they used a text retriever and image based retrieval using RIS, but they provide the retrieved sources directly to MLLM such as GPT-5.1 before generating a question. With a few shot examples, the MLLM fromulates 10 questions and answers based on the retrieved evidence, and then uses these question answer pairs to predict the veracity of the claim. This approach is more efficient as it does not require a separate question generation step, but it may be less interpretable as the MLLM is generating both the questions and answers based on the retrieved evidence.
Here question answer pairs are generated in a single step after retrieving the evidence, which may not allow for a detailed analysis of the evidence or the claim, and may also lead to less interpretable predictions.
Furthermore, the approach relies heavily on the quality of the retrieved evidence, and if the retrieval step fails to retrieve relevant information, the question-answer generation and subsequent veracity prediction may be compromised.
Also there is risk of hallucination in the question-answer generation step, where the MLLM may generate questions and answers that are not grounded in the retrieved evidence, leading to inaccurate veracity predictions.

~\citet{Jung2026Mar} VILLAIN follow similar RAG based approach, using text retriever and visual retriever to retrieve evidence, but they have added a extra step where three specialized VLM based agents generate a detaield analysis reports focused on  a specific evidentiary relationship. A Text-Text agent to analyze the alignment between the claim text and the retrieved textual sources.
A Image-Text agent to Examines the relationship between the claim images and the textual evidence derived from them and a cross modal agent Investigates the complex interactions between all textual and visual evidence sets to identify global narratives or critical discrepancies. The analysis reports are synthesized by a Lead Fact-Checking Adjudicator agent into 20 diagnostic QA pairs, 5 pairs in 4 iterations. The verdict prediction egent is predicted by filtering the 20 QA pairs to 10 QA pairs and generate a verdict and justification.
This approach allows for a more detailed analysis of the evidence and the claim, and may lead to more interpretable predictions. However, it is also more complex and may require more computational resources, which can limit its practicality for real-world deployment.
It is retrieval dependent and there is difficulty in scaling such systems.


\begin{table}[ht]
\centering
\caption{Comparison of top-performing multimodal fact-checking systems in the AVeriMaTeC Shared Task.}
\label{tab:comparison_systems}
\begin{tabular}{l c c c}
\hline
\textbf{Feature} & \textbf{AIC CTU} & \textbf{ADA-AGGR} & \textbf{VILLAIN} \\
\hline
Rank & 3rd Place & 2nd Place & 1st Place \\
MLLM Backbone & GPT-5.1 & Gemini 3 Pro & Gemini-2.5-Pro \\
Complexity & 1 MLLM call & Sequential (3+ steps) & 8 MLLM calls (Multi-agent) \\
Key Innovation & Low-cost modular RAG & Image downsampling for speed & URL filling \& Multi-agent analysis \\
Primary Goal & Accessibility/Cost & Deployment/Efficiency & High-Veracity Performance \\
\hline
\end{tabular}
\end{table}

Retrieved evidences are the backbone of any fact checking system. Robust multimodal retrieval remains a major challenge in multimodal fact-checking, as the quality of retrieved evidence directly impacts the accuracy of verification and in some cases the resukting QA pairs and the final verdict may be more dependent on the quality of the retrieved evidence than the reasoning capabilities of the model, which highlights the importance of improving retrieval techniques for multimodal fact-checking systems.
In this competition with a already set knowledge store, retrieving evidences to generate QA pairs is fine but in a real world setting, the system may need to counter question the claim iteratively to retrieve evidences from web.
While these recent state-of-the-art frameworks like VILLAIN and AIC CTU rely on the mxbai embedding family and standard reverse image search, my approach introduces a more nuanced retrieval ensemble. By combining SPLADE for learned sparse expansion, BGE-en for dense textual semantics, and LongCLIP for long-context visual grounding, I address the 'vocabulary mismatch' and 'description bottleneck' limitations inherent in current multimodal RAG pipelines

\section{Retrieval Systems for Fact-Checking}


Sparsness is needed primarily for efficiency in large scale retrieval systems. Sparse retrieval is a technique in information retrieval that finds relevant documents by matching exact words or tokens between a query and documents. It creates high-dimensional, mostly zero-valued vectors for text,
which allows for highly efficient search via inverted indices. For a retrieval system to be efficient, it must avoid calculating scores for every document in a collection. Instead, it should only consider documents that contain at least one query term. 

Sparse retrieval identifies relevant documents based on statistical word frequency information. It is "sparse" because it typically uses an inverted index where most document-term values are zero; a document is only indexed or "scored" if it explicitly contains a query term.
Sparse retrieval is required due to the following reasons:
\begin{itemize}
    \item \textbf{Efficiency with Large Datasets}: The system must ``quickly obtain'' relevant documents from a collection of millions of open-domain Wikipedia documents. Sparse methods are computationally efficient enough to handle this scale.
    \item \textbf{Maximizing Recall}: The primary goal of this initial sparse stage is to achieve ``as high as possible recall''. This ensures that the potential evidence is captured in a smaller candidate pool before more complex, computationally intensive models are used for refinement ~\citet{Wang2022Aug}.
\end{itemize}

In fact-checking tasks, sparse retrieval remains particularly effective because many verification problems depend on exact entity matching, numerical consistency, and precise lexical evidence. Claims involving names, dates, locations, or quoted statements often require retrieval systems to prioritize exact term overlap rather than semantic similarity alone.


TF-IDF (Term Frequency–Inverse Document Frequency) can be seen as the most basic form of sparse reteival.
TF-IDF (Term Frequency-Inverse Document Frequency) is a traditional information retrieval model that scores documents based on the product of two components: a weight derived from the frequency of a term in a document (tf) and a weight based on how many documents in the collection contain that term (idf).


\begin{itemize}
    \item \textbf{Term Frequency (tf)}: This component represents how many times a query term appears in a document. In traditional tf$\times$idf models, the tf component is often linear, meaning the weight increases proportionally with each additional occurrence of the term.
    \item \textbf{Inverse Document Frequency (idf)}: This component weights terms based on their rarity across the entire collection. The sources define a specific idf weighting formula ($w_{\text{IDF}_i}$) that is used in the absence of relevance information, which serves as a ``close approximation to classical idf''.
\end{itemize}

BM25 improves upon traditional TF-IDF (Term Frequency-Inverse Document Frequency) by introducing a more nuanced, non-linear approach to term weighting and document length ~\citet{Robertson2009Sep}. 
BM25 is a long-standing standard in information retrieval primarily valued for its effectiveness and simplicity
. It operates on a probabilistic framework that ranks documents by analyzing term frequency and inverse document frequency to capture relevance
. One of its primary advantages is its computational efficiency, as it handles inputs—particularly short queries—with minimal processing overhead and significantly lower latency.
It is more effective at matching salient, rare phrases or specific names
. For instance, BM25 can better retrieve a passage containing a unique name like "Thoros of Myr," which a dense encoder might lack the capacity to represent specifically ~\citet{Karpukhin2020Nov}.
BM25 was used for text retrieval for the claim text in the baselines of the Fever competition ~\citet{Cao2025May} and also by ~\citet{Upravitelev2026Mar}.

The primary limitation of the BM25 model is its reliance on keyword matching, which hinders its ability to handle complex and semantically rich queries. While it has been a long-standing standard in information retrieval due to its simplicity, it often struggles to capture the deeper semantic meaning of documents and queries beyond simple term overlap. Consequently, in large-scale retrieval scenarios involving billions of documents, BM25 consistently demonstrates lower retrieval effectiveness—measured by metrics like MRR@10 and Semantic Similarity Scores (SSS)—compared to modern neural sparse models.
SPLADE addresses these limitations by utilizing learned sparse lexical representations that allow for superior semantic understanding and enhanced query expansion. By incorporating learned embeddings, SPLADE captures deeper semantic relationships that traditional keyword-based models miss, leading to significant improvements in retrieval performance, particularly for complex inputs.
However splade models incure significantly higher FLOPS and latency compared to traditional BM25 models struggles with efficicnecy when handling shot inputs.
also splade is limited by the BERT Wordpiece vocabulary which can constrain their overall lexical representation power and flexibility in document-query interactions.
SPLADE, while a significant improvement over traditional sparse methods, still relies on term expansion to capture meaning.~\citet{Won2025Nov}.

Dense Passage Retrieval (DPR) represents a shift in retrieval methods from traditional sparse, keyword-based models like BM25 and splade to a dense dual-encoder framework
. While BM25 and splade relies on high-dimensional sparse vectors, exact token matching and lexical search, DPR utilizes BERT-based encoders to map questions and passages into a low-dimensional continuous space, measuring relevance via dot-product similarity
.This dense representation allows the system to capture latent semantic relationships and lexical variations (e.g., equating "bad guy" with "villain") that keyword-matching systems often fail to retrieve. Empirically, DPR significantly outperforms BM25 on most benchmarks, improving top-20 retrieval accuracy by 9\% to 19\% absolute ~\citet{Karpukhin2020Nov}.
~\citet{Ullrich2026Mar} and ~\citet{Jung2026Mar} both have use dense embedding models for text retreival with models mxbai-embed-large-v1 and mxbai-embed respectively.

\section{Hybrid Retrieval}

While they are highly effective at semantic matching, they have specific limitations. Dense vectors are designed for context but lack the specific keyword information required for exact matching. Because they focus on semantic meaning rather than specific lexical tokens, they may struggle in scenarios where a specific name, code, or technical term is critical to the query's relevance.
 show that a deep
retrieval model poorly generalizes to a new domain with large domain shift, while
lexical matching and expansion models are robust across domains ~\citet{Chen2022Jan}.
Classic lexical retrieval models struggle to understand the underlying meanings
of queries and documents. Neural embedding based retrieval models can soft
match queries and documents, but they lose specific word-level matching information ~\citet{Gao2020Apr}.

A hybrid strategy, which combines dense vectors with sparse vectors (such as BM25 or SPLADE), is used to exploit the complementary strengths of both methods while mitigating their individual weaknesses.

\begin{itemize}
    \item \textbf{Complementary Strengths}: Sparse methods excel at exact keyword matching and provide interpretable representations, filling the gap left by dense retrievers.
    \item \textbf{Enhanced Accuracy}: Empirical evidence shows that fusing these two representations significantly improves overall retrieval accuracy compared to using either method alone ~\citet{Gao2020Apr} ~\citet{Chen2022Jan}.
    \item \textbf{Robustness in Specific Scenarios}: Hybrid strategies demonstrate improved performance and robustness in cross-domain retrieval (where the search engine is used on a different domain than its training data) and long-text retrieval ~\citet{Luan2021}.
    \item \textbf{Support for Modern AI}: Hybrid retrieval has become a standard practice in Retrieval Augmented Generation (RAG) for Large Language Models (LLMs) to ensure the most relevant context is provided to the model.
\end{itemize}

However combining these two methods can be challenging. Dense and sparse vectors have vastly different distributions.
Dense vectors are typically normalized (ranging 0 to 1), whereas sparse vectors have indeterminate magnitudes and different similarity probability spaces
This makes it difficult to find an optimal fusion weight (α) to balance the two types of distances accurately.

ost prior works [42,17,24,23] model this combination as a linear interpolation of the scores of deep and lexical retrieval models. This fusion method is
sensitive to the score scales and the weights assigned to the different models [42],
which needs careful score normalization and weight tuning, especially when multiple models are combined. We expect that the raw scores of the models can vary
from one domain/dataset to another, and likewise the interpolation weights ~\citet{Wang2021Jul}.

There are approaches where the two retrievals are done separately from their own indexes and then a Reciprocal rank fusion
is done to generate the final ranking by considering the ranking positions of each candidate generated by different
models, instead of fusing their scores. RRF demonstrates robust and effective
ensembles. It can be easily applied to a
new domain in a zero shot setting (with no in-domain training data), we would
like to eliminate such domain-specific normalization and tuning ~\citet{Chen2022Jan}.

Rather than using two retreivals from different indices and combining them with RRf ~\citet{Zhang2024Oct} discuss about creating a unified inde to handler hybrid vectors.
This typically involves concatenating the dense and sparse vectors into a single high-dimensional representation. A single search algorithm, such as Hierarchical Navigable Small World (HNSW) or an Inverted File index (IVF), is then used to retrieve results from this unified space.
However it has its downsides, as the Dense and sparse vectors have vastly different distributions and since Sparse vectors are high-dimensional, and calculating their distances is significantly more expensive than dense calculations, In a unified graph-based index, sparse distance computation can account for over 60\% of total search time and take approximately 11 times longer than dense distance calculations.

\section{Multimodal Retrieval}

CLIP, which stands for Contrastive Language-Image Pre-training, is a computer vision system developed to learn state-of-the-art image representations directly from natural language supervision.
The system consists of an image encoder and a text encoder that are trained together and the feature respresentations form each encoder are linearly projected into a shared embedding space.Alignment is achieved by maximizing the cosine similarity between the embeddings of the correct image-text pairs while minimizing the similarity for all incorrect pairings in the batch
. This process "connects" visual representations to language, allowing the model to understand visual concepts through the descriptions used to reference them.
Because it understands the relationship between text and images, CLIP is highly effective at finding relevant images in a database based on a text query, or finding relevant text given an image.
The model has learned to recognize digitally rendered text, enabling it to perform tasks like reading signs or identifying content in memes.
CLIP excels at recognizing verbs and activities (e.g., "hammering a phone" or "drilling an egg"). Because CLIP learns such a broad set of concepts from 400 million internet pairs, it displays exceptional zero-shot retrieval performance ~\citet{Radford2021Feb}.

CLIP imposes a hard ceiling of 77 tokens for input text due to its absolute positional embedding. Because CLIP was trained predominantly on brief summary texts, its image encoder tends to focus only on the most crucial elements, disregarding finer details and relationships within an image.
Long-CLIP improves the model by unlocking long-text capabilities while preserving the original strengths of CLIP.
To overcome the 77-token limit, Long-CLIP extends the maximum input length to 248 tokens. It retains the first 20 positional embeddings, which were already well-trained in the original model.
It then interpolates the remaining 57 embeddings using a larger ratio
. This minimizes disruption to the established representation of short text positions while allowing the model to process much longer sequences.
Long-CLIP shows a 25\% improvement in long-caption image-text retrieval and a 6\% improvement in traditional short-text retrieval tasks like COCO and Flickr30k

~\citet{Samuel2025Jul} have used SigLIP(Sigmoid Loss for Language-Image Pre-training), to bridge the modality gap between natural language text queries and visual video frames.
At search time, the SigLIP text encoder processes the user's query, which is then matched against the frame embeddings to find the best visual match for each video.
SigLIP is often preferred over the original CLIP because it replaces standard softmax normalization with a sigmoid loss function, which allows for better scalability and efficiency during training.
They achieved an 81\% improvement in nDCG@20 over leading multimodal encoders and a 37\% improvement over single-modality systems on the MultiVENT 2.0 dataset.

In the MMMORRF system from ~\citet{Samuel2025Jul}, Reciprocal Rank Fusion (RRF)—specifically a novel variant called Weighted Reciprocal Rank Fusion (WRRF)—is used to combine the rankings from separate retrieval sub-pipelines (vision-based and text-based) into a single, unified result.
WRRF allows the system to capitalize on the complementary strengths of different modalities, such as visual frames (processed by SigLIP), spoken audio (ASR), and embedded text (OCR).
Unlike traditional RRF, which treats all sources equally, WRRF uses video-dependent importance weights
. This means the system can adaptively prioritize the modality that is most likely to be reliable for a specific piece of content.

Similarly in ~\citet{Mourao2014Jun}, Multimodal retrieval is typically implemented as a late fusion task, where information from different modalities (such as text, images, or audio) is indexed and searched separately before being combined into a single, unified result list.
Each independent search generates its own ranked document list, where items are assigned individual scores and ranks based on that specific modality's relevance.
Along with Reciprocal Rank (RR) and Reciprocal Rank Fusion (RRF), A more advanced variant, Inverse Square Rank (ISR), uses a quadratic decay (1/rank 
2) to prioritize top-ranked results and multiplies the final score by the frequency of the document across lists to exploit the "Chorus Effect".
The chorus effect occurs when multiple different retrieval approaches all retrieve the same item, suggesting a higher probability that the item is relevant to the query.
Methods such as CombSUM, CombMAX, and CombMNZ were also used by aggregating the actual relevance scores from each engine, however these can be unstable because score distributions vary significantly between different search systems.

\section{Evidence Ranking and Diversification}

Cross-encoders are primarily utilized for reranking because they provide significantly higher accuracy and better generalization than bi-encoders like splade and dense retrievers, particularly in zero-shot or out-of-domain scenarios.
. While they are computationally more expensive, they are "necessary" in a multi-stage retrieval pipeline to refine the results provided by faster, less precise first-stage retrievers like BM25,splade and dense retreivals.
Unlike bi-encoders, which represent queries and documents as separate vectors, cross-encoders process the query and document together
. This allows for "early query-document interactions" within the model's architecture, which is a significant factor in their superior effectiveness.
Cross-encoders demonstrate a much stronger ability to generalize to new domains that were not seen during fine-tuning
. While bi-encoders often perform similarly to cross-encoders on in-domain tasks, their performance tends to drop significantly—or "extrapolate substantially worse"—when applied to new domains ~\citet{Rosa2022Dec}.
~\citet{Rosa2022Dec} uses cross encoders like monoT5 as rerankers in the second stage of the pipeline and is taksed with sorting a candidiate list of 1000 texts that have be retrieved by BM25 and 
concludes that the reranker's effectiveness remains consistently high regardless of whether the initial 1,000 documents were provided by a sparse retriever (BM25) or a dense retriever (GTR).
Similarly ~\citet{Ma2023May} uses LLMs like GPT-3 as rerankers  where it takes query and a list of multiple documents simultaneously and generate a reordered list of document identifiers based on their relevance.

However, while cross-encoders are effective in reranking, a lot of duplicate or same evidences can score high and those homogenous documents could be high in the ranks which will not involve other alternative voices in the evidences.
MMR - Maximal Marginal Relevance is an innovative method used to enhance the performance of Retrieval-Augmented Generation (RAG) pipelines in automated fact-checking.
MMR is used for diversity-based reranking when reordering documents
. This ensures that the system doesn't just retrieve the most similar documents, which might all contain the same information, but instead selects a diverse set of evidence.
By reordering documents and producing summaries that prioritize non-redundant information, systems can provide a more comprehensive set of facts to the reasoning model.

According to ~\citet{Carbonell1998Aug}, in scenarios where there is a "vast sea" of potentially relevant documents, many often contain duplicative or overlapping information
. MMR identifies "marginal relevance," meaning a document only ranks highly if it provides new information that hasn't appeared in previously selected results.
The method is user-tunable through a parameter (λ) that balances relevance and diversity
. Users can choose a smaller λ to get a broad, panoramic sampling of a topic or a larger λ to focus on the most reinforcing relevant documents ~\citet{Ullrich2024Nov}.


Similar to MMR, ~\citet{Do2025Apr} have use a Diverse Length-aware MMR(DL-MMR) which ensures that retrieved exemplars have diverse target lengths
,
. This is used because Large Language Models (LLMs) tend to generate summaries that match the length of the examples provided to them; providing diverse lengths prevents the LLM from being "locked" into a specific length.
A primary reason for using the DL-MMR variant is its extreme efficiency compared to standard MMR
,
. Standard MMR requires comparing every pair of exemplars in a pool, which is computationally expensive (n(n−1)/2 comparisons), whereas DL-MMR only needs to store and compare fixed length values (n comparisons).
Using MMR variants like DL-MMR can reduce memory usage by over 781,000 times and computational costs by over 500,000 times while maintaining the same level of informativeness in the final generated summary.


\section{LLM and MLLM for Fact-Checking}

Other important part of a fact checking system is the reasoning module, which is responsible for analyzing the retrieved evidence and making a final prediction about the veracity of the claim. 


Large Language Models (LLMs) and Multimodal Large Language Models (MLLMs) have become increasingly popular for this task due to their ability to understand and reason over complex information.
LLMs can evaluate claims based on plausibility, language formality, and narrative structure
. This allows them to perform surprisingly well even on claims about events that occurred after their training data cut-off.
Research shows LLMs are particularly effective at identifying and classifying claims expressed as opinions, as they can rely on tone and style rather than just external world knowledge.
Despite these uses, LLMs are not infallible; they often struggle more with factual-sounding claims than opinions and show significant performance variability across different languages and topics ~\citet{Tang2024Nov}.

LLMs are used to decompose claims into questions, draw out context and meaning from the evidences and reason over the retrieved evidences to make a final prediction about the veracity of the claim. 
~\citet{Schlichtkrull2023May}, ~\citet{Schlichtkrull2024Nov},~\citet{Ullrich2024Nov}, ~\citet{Ullrich2025Aug} have all relied on LLMs to generate question-answer pairs based on the retrieved evidence, and then use these question-answer pairs to make a final prediction about the veracity of the claim.

After the introduction of the multimodal AVERIMATEC dataset, ~\citet{Cao2025May} , most the competitive system have employed MLLMs to explore the multimodal nature of the claims and reasoning over the retrieved multimodal evidences.

\section{QA generation in fact checking}

Question-Answer (QA) pairs are useful for fact-checking because they break down complex claims into manageable parts, facilitating the retrieval of evidence and improving the efficiency and clarity of the verification process.

~\citet{Setty2024Jul} have shown that question decomposition from claim for fact checking can be effectively automated and demonstrate that smaller generative
models, fine-tuned for the question generation task using data augmentation from various datasets, outperform large language models
by up to 8\%. They also show the evidence retrieved
using machine-generated questions proves to be significantly more
effective for fact-checking than that obtained from human-written
questions.

In automated systems like Varifocal ~\citet{Ousidhoum2022Dec}, this is done by first extracting "focal points"—coherent phrases from the claim and its metadata (source and date)—using a dependency parser
. A model then generates unique questions for each focal point to scrutinize the information, followed by a re-ranking process that removes duplicates and selects the most relevant and informative questions for investigation.

The baseline paper for Averimatec ~\citet{Cao2025May} explores three question-generation strategies for fact checking: Parallel Question Generation (PQG), which generates all questions at once; Dynamic Question Generation (DQG), where each question depends on prior evidence and earlier questions; and Hybrid Question Generation (HQA), which combines both by generating initial questions in parallel and later ones dynamically.

~\citet{Jung2026Mar}, ~\citet{Ullrich2025Aug} and ~\citet{Schlichtkrull2023May} have employed in-context learning with retreived several few-shot QA-pair generation examples from the
training data using BM25, and these examples are
incorporated into the QA generation prompt.

\section{Evaluation of Fact-Checking Systems}

Strong evaluation techniques are required to ensure that fact-checking systems are not only accurate but also robust and generalizable across different domains and claim types.
Early evaluation methods in FEVER Workshop 2021 ~\citet{Aly2021Nov}, focused on  a single score.For any given claim, a system receives a score of 1 only if both of the following conditions are met- The predicted verdict label is correct and
The predicted evidence contains at least one complete gold evidence set (meaning the predicted evidence is a superset of at least one required evidence set provided by human annotators). exact matching is done.

Later both ~\citet{Schlichtkrull2023May} and ~\citet{Schlichtkrull2024Nov} used a scoring function called Hungarian METEOR to measure how well a system's generated questions and answers match reference sets, as the the same evidence can be found in many different sources or phrased in various ways. which relies on
approimate matching rather than absolute matching. To ensure veracity predictions are grounded in actual facts, a system only receives credit if its evidence meets a quality threshold of 0.25 on the Hungarian METEOR scale.

Recently in FEVER 9  ~\citet{Cao2025May} competition, with the introuduction of multimodal eveidences, These are compared against human-annotated "gold" references using LLM-based scoring (Gemini-2.0-Flash) in a two-stage process that checks both textual and visual similarity. The quality of the generated justifications is measured by comparing them to human ground truth using ROUGE-1 and Ev2R, a metric shown to align well with human assessment for open-ended generation tasks.



































